\clearpage\selectlanguage{german}\fancyhead[L]{\small BEMERKUNGEN.}\pdfbookmark[0]{Bemerkungen}{nachlass-bemerkungen}\SourceLeaf{239}
\begin{center}\textbf{BEMERKUNGEN.}\par\smallskip\rule{.15\linewidth}{.4pt}\end{center}
In der Handschrift bildet die \textit{Determinatio seriei nostrae per aequationem differentialem etc.} den zweiten Teil einer Abhandlung mit dem Titel \lq\textit{Disquisitiones generales circa functiones a serie infinita $1+\frac{\alpha.\beta}{1.\gamma}x+$ etc. pendentes auctore} \textsc{Carolo Friderico Gauss} \textit{societati regiae traditae Nov.} 1811.\rq\ Der erste Teil ist nach einer ausführlicheren Ausarbeitung der Einzelheiten in die bekannte, 1812 am 30. Januar veröffentlichte Abhandlung übergegangen. Dabei wurden aus 27 Artikeln 37, und die 65 nummerierten Gleichungen wurden auf 79 vermehrt. An diese schließen sich in der Handschrift die Nummern der Artikel und Gleichungen des zweiten Teils an; zur Bequemlichkeit des Lesers hat man sie jedoch auf die Nummern der veröffentlichten Abhandlung folgen lassen und auch die vorkommenden Verweise auf diese bezogen.

Das Handexemplar der gedruckten Abhandlung \textit{Disqu. gen. circa seriem etc.} enthält zu Artikel 28 die Aufzeichnung
\[
\begin{aligned}&\textit{Die beste Definition von }\Pi\textit{ ist, dass}\\&
\Pi m=\int_{-\infty}^{+\infty}e^{(m+1)x}e^{-e^x}\,\mathrm dx\end{aligned}
\]
ferner die Gleichung 58 auch noch in dieser Form:
\[
\begin{aligned}&[58]\qquad \log\Pi z=(z+\tfrac12)\log z-z+\tfrac12\log2\pi\\&+\tfrac16.\tfrac12 P+\tfrac18.\tfrac23 Q+\tfrac1{10}.\tfrac34 R\\&+\tfrac1{12}.\tfrac45 S+\text{usw.}\end{aligned}
\]
Sie lässt sich unmittelbar aus 61, 62 und 57 herleiten.

In einem Notizbuch ist mithilfe der Gleichung 58 in der \textsc{Euler}'schen Form der Wert von $\log\Pi(10+i)$ und daraus der von $\log\Pi i$ als
$9{,}7173075-17^\circ16'57''693\,i$ und
$\Pi i=+0{,}4980156-0{,}1549496i$ berechnet (nach dem Jahr 1847).

Artikel 48 ist hier so wiedergegeben, wie er nach vielfachen Streichungen von Wörtern und ganzen Sätzen in der Handschrift gelesen werden muss; nach der Absicht des Verfassers dürften jedoch wohl noch die Wörter \lq\textit{e solo theoremate binomiali}\rq\ wegzulassen sein.
\begin{flushright}\textsc{Schering.}\end{flushright}
\begin{center}\rule{.4\linewidth}{.4pt}\end{center}
